%! TEX root = ../main.tex

\clearpage

\section{Các tính năng của hệ thống}
\textit{<Tài liệu này minh họa việc tổ chức các yêu cầu chức năng cho sản phẩm theo các tính năng của 
hệ thống, các dịch vụ chính được cung cấp bởi sản phẩm. Ta có thể tổ chức mục này theo use 
case, chế độ vận hành, lớp người sử dụng, lớp đối tượng, sự phân cấp theo chức năng, hoặc kết 
hợp chúng lại.>}

\subsection{Tính năng 1 của hệ thống}

\textit{<Không nói là “Tính năng 1 của hệ thống” mà viết cụ thể tên chức năng.>}
 
\subsubsection{Mô tả và mức ưu tiên}
\textit{<Cung cấp một mô tả ngắn gọn về tính năng và chỉ ra mức độ ưu tiên của nó (cao, 
trung bình, hay thấp). Ta cũng có thể đưa vào cách xếp mức ưu tiên dựa vào các 
tiêu chí như lợi ích, chi phí, rủi ro (mỗi cái được xếp loại từ thấp – 1 đến cao - 9).>}

\subsubsection{Tác nhân / Chuỗi đáp ứng}
\textit{<Liệt kế các chuỗi hoạt động của người sử dụng và các đáp ứng của hệ thống mà 
chúng kích thích hành vi được định nghĩa bởi tính năng này. Chúng tương ứng với 
các thành phần của dialog được liên kết với các use case.>}

\subsubsection{Các yêu cầu chức năng}
\textit{<Ghi thành từng nhóm các yêu cầu chức năng chi tiết có liên quan tới tính năng này. 
Chúng là các khả năng của phần mềm và phải có để người sử dụng thực hiện các 
dịch vụ được cung cấp bởi tính năng này, hay thực thi use case. Ta cũng nên trình 
bày cách thức sản phẩm đáp lại các điều kiện lỗi được dự đoán trước hay các input 
không hợp lệ. Các yêu cầu nên ngắn gọn, hoàn chỉnh, không mơ hồ, có thể kiểm tra 
và cần thiết. Sử dụng “TBD” khi thông tin cần thiết vẫn chưa có sẵn.> }

\textit{<Mỗi yêu cầu nên được đặt tên duy nhất bằng một số tuần tự hay một thẻ có 
nghĩa.>}

REQ-1: \\
\indent REQ-2:

\subsection{Tính năng thứ 2 của hệ thống (và cứ như thế)}